% 前言 (原书 Introduction, 第 ix 页)
\chapter*{前言}
\markboth{前言}{前言}
\addcontentsline{toc}{chapter}{前言}

\begin{center}
\itshape 献给 John——他始终不离不弃, 陪伴我们历经变迁.
\end{center}

\vspace{2em}

本书旨在对 Bochner--Riesz 平均的几何理论作一个自足的阐述. 这一主题处理 Fourier 分析中最基本的问题: Fourier 级数何时收敛到原来的函数? 二十世纪七十年代中期, 这一问题取得了实质性的进展, 但相关技巧至今仍只能见于专业文献. 我们的意图是给出一个研究生能够读懂的论述. 为了保持叙述的一致性, 我们对某些重要课题作了舍弃. 我们假定读者熟悉研究生水平的实分析, 并熟悉分布理论与 Fourier 变换的基本事实. 一本基本的参考书是 Stein 与 Weiss 的 \emph{Introduction to Fourier Analysis on Euclidean Spaces}~\cite{steinweiss}, 以及 Rudin 的两部著作\footnote{译者注: 原书此处引用编号重复作 [43] 与 [43], 应为 [42] 与 [43].}~\cite{rudinrca,rudinfa}.

在写作本书的过程中, 多年来我们从与同事们的广泛交谈中获益匪浅. 我们感谢 E.~Fabes, R.~Fefferman 与 E.~M.~Stein 教授对第~1、2、3 章材料的帮助. 第~4、5 章的内容受到了与 A.~W.~Knapp 教授交谈的影响. 对于第~7、8 章的一般思想, 我们感激 A.~Córdoba 与 C.~Fefferman 教授. 本书初稿写作期间, 第一作者得到了美国国家科学基金会资助 (NSF grants MCS 8202165 与 8001799). 第一作者感谢 W.~Beiglböck 教授与海德堡大学应用数学研究所, 使他有机会完成本书的倒数第二稿. 第二作者感谢中华人民共和国教育部在整个项目期间给予的支持. 两位作者都感谢 K.~Merryfield 教授对本书早期版本所作的细致审读.

两位作者感谢 University Research Institute 在本书成书过程中给予的经费支持.

最后, 第一作者感谢 Margaret Combs 与 Jan Duffy, 她们总是乐于提供帮助.
